\begin{recipebox}{Mince \& Cheese Pie Bake}{4}{15 min prep + 40 min}

A Kiwi mince and cheese pie in a baking dish --- no pie tins or bases needed.
Roughly 35\,g protein per serve; serve with green veg.

\begin{ingredients}
  \item[] \textbf{Filling}
  \item 500\,g lean beef mince
  \item 1 onion, diced
  \item 2 carrots, grated
  \item \nicefrac{1}{4} cup dry red lentils, rinsed
  \item 1\nicefrac{1}{2} cups beef stock
  \item 2 tbsp tomato paste
  \item 1 tbsp Worcestershire sauce
  \item 1 tbsp cornflour mixed with 2 tbsp cold water
  \item 1 cup frozen peas
  \item Salt and pepper
  \item[] \textbf{Top}
  \item 75\,g grated tasty cheese
  \item 1 sheet puff pastry, thawed (or 6 sheets filo brushed with a little oil)
  \item 1 egg, beaten, or milk to brush
\end{ingredients}

\begin{method}
  \item Preheat oven to 200\degree C bake.
  \item Brown the mince and onion in a frying pan. Add the carrot, lentils, stock, tomato paste and Worcestershire and simmer for 15 minutes, until the lentils are soft.
  \item Stir in the cornflour mixture and peas and simmer until thick. Season. Tip into a baking dish (about 20 $\times$ 25\,cm) and let it cool for 5 minutes.
  \item Scatter the cheese over the filling. Lay the pastry over the top, trim and press the edges, and cut a couple of slits. (For filo, scrunch the oiled sheets loosely on top.)
  \item Brush with egg and bake for 20--25 minutes, until the pastry is puffed and golden.
\end{method}

\begin{notelist}
  \item Filo is lighter and lower in saturated fat than puff pastry.
  \item Letting the filling cool a little before adding the pastry stops it going soggy.
\end{notelist}

\end{recipebox}
